\chapter{Path Integral for Scalar Fields}

\begin{remarkbox}
The path integral for a particle sums over histories \(q(t)\). The path integral
for a field sums over field configurations \(\phi(x)\). This chapter develops the
functional integral for scalar fields and shows how propagators, correlation
functions, perturbation theory, and diagrams arise from a single generating
functional.
\end{remarkbox}

\section{From Paths to Field Configurations}

In ordinary quantum mechanics, the path integral is
\[
    \int\mathcal{D}q(t)\,e^{iS[q]}.
\]
For a scalar field, the dynamical variable is not a single function of time but
a function of spacetime. The corresponding path integral is
\[
    \int\mathcal{D}\phi\,e^{iS[\phi]}.
\]
The measure \(\mathcal{D}\phi\) symbolically means integration over all field
configurations.

For a real scalar field, the action is
\[
    S[\phi]
    = \int\dd^4x\,
      \left(
          \frac{1}{2}\partial_\mu\phi\partial^\mu\phi
          -\frac{1}{2}m^2\phi^2
      \right)
\]
for the free theory. Interactions are added by including non-quadratic terms,
for example
\[
    S[\phi]
    = \int\dd^4x\,
      \left(
          \frac{1}{2}\partial_\mu\phi\partial^\mu\phi
          -\frac{1}{2}m^2\phi^2
          -\frac{\lambda}{4!}\phi^4
      \right).
\]

\section{The Generating Functional}

Introduce a source \(J(x)\) coupled linearly to the field:
\[
    Z[J]
    = \int\mathcal{D}\phi\,
      \exp\left(iS[\phi]+i\int\dd^4x\,J(x)\phi(x)\right).
\]
The source is an auxiliary function. It is used to generate correlation
functions and then set to zero.

The vacuum-to-vacuum functional is
\[
    Z[0]=\int\mathcal{D}\phi\,e^{iS[\phi]}.
\]
Normalized correlation functions are obtained by dividing by \(Z[0]\). For
example,
\[
    \langle0|T\hat\phi(x)\hat\phi(y)|0\rangle
    = \frac{1}{Z[0]}
      \left.\frac{1}{i^2}
      \frac{\delta^2 Z[J]}{\delta J(x)\delta J(y)}\right|_{J=0}.
\]
Thus the path integral generates the same time-ordered correlators that appeared
in the operator formulation.

\section{The Free Scalar Generating Functional}

The free scalar action can be written, after integrating by parts and ignoring
boundary terms, as
\[
    S_0[\phi]
    = -\frac{1}{2}\int\dd^4x\,\phi(x)(\Box+m^2)\phi(x).
\]
With a source, the exponent is quadratic in \(\phi\):
\[
    iS_0[\phi]+i\int\dd^4x\,J\phi.
\]
This is an infinite-dimensional Gaussian integral. The result has the form
\[
    Z_0[J]
    = Z_0[0]\exp\left(
        \frac{i}{2}\int\dd^4x\dd^4y\,
        J(x)D_F(x-y)J(y)
      \right),
\]
where \(D_F\) is the Feynman propagator satisfying
\[
    (\Box_x+m^2)D_F(x-y)=-i\delta^{(4)}(x-y).
\]
The inverse of the quadratic operator in the action is therefore the propagator.
This is the path-integral version of the result obtained from canonical
quantization.

\section{Sources and Functional Derivatives}

Functional derivatives insert fields. The basic rule is
\[
    \frac{1}{i}\frac{\delta}{\delta J(x)}Z[J]
    = \int\mathcal{D}\phi\,\phi(x)
      \exp\left(iS[\phi]+i\int\dd^4z\,J(z)\phi(z)\right).
\]
Therefore
\[
    \langle T\phi(x_1)\cdots\phi(x_n)\rangle
    = \frac{1}{Z[0]}
      \left.\frac{1}{i^n}
      \frac{\delta^n Z[J]}{\delta J(x_1)\cdots\delta J(x_n)}\right|_{J=0}.
\]
For the free theory, differentiating \(Z_0[J]\) reproduces Wick's theorem. The
exponential quadratic in \(J\) generates all pairings of the external points,
and each pair contributes one propagator.

\section{Gaussian Functional Integrals}

The free scalar generating functional is the infinite-dimensional analogue of
\[
    \int \dd^n x\,
    \exp\left[-\frac{1}{2}x^TAx+J^Tx\right]
    \propto
    \exp\left(\frac{1}{2}J^TA^{-1}J\right).
\]
In field theory, the matrix \(A\) becomes a differential operator, and
\(A^{-1}\) becomes a Green function.

Schematically,
\[
    A = \Box+m^2,
    \qquad
    A^{-1} \sim D_F.
\]
The \(i\epsilon\) prescription specifies which inverse is meant. Without it, the
operator has several possible Green functions. With it, the inverse is the
Feynman propagator appropriate for time-ordered vacuum correlators.

The determinant of the quadratic operator contributes to \(Z_0[0]\). In many
normalized correlation functions this factor cancels, but determinants become
important in loop effects, effective actions, and vacuum energies.

\section{Interacting Scalar Theory}

For \(\phi^4\) theory,
\[
    S[\phi]
    = S_0[\phi]-\int\dd^4x\,\frac{\lambda}{4!}\phi(x)^4.
\]
The generating functional is
\[
    Z[J]
    = \int\mathcal{D}\phi\,
      \exp\left(iS_0[\phi]-i\int\dd^4x\,\frac{\lambda}{4!}\phi^4
      +i\int\dd^4x\,J\phi\right).
\]
The interaction can be represented as functional derivatives acting on the free
generating functional:
\[
    Z[J]
    = \exp\left[-i\int\dd^4x\,\frac{\lambda}{4!}
      \left(\frac{1}{i}\frac{\delta}{\delta J(x)}\right)^4\right]Z_0[J].
\]
Expanding the exponential in powers of \(\lambda\) gives perturbation theory.
This is the path-integral version of the Dyson expansion.

\section{Diagrammatic Expansion}

The diagrammatic rules follow from expanding the interacting generating
functional and differentiating with respect to sources. Each factor of
\(D_F(x-y)\) becomes a line. Each factor of \(-i\lambda\int\dd^4x\) becomes an
interaction vertex integrated over spacetime.

For \(\phi^4\) theory:
\begin{itemize}
    \item a propagator contributes \(D_F(x-y)\) in position space or
    \(i/(k^2-m^2+i\epsilon)\) in momentum space;
    \item a vertex contributes a factor \(-i\lambda\);
    \item four scalar lines meet at each vertex;
    \item one integrates over internal vertex positions or internal momenta;
    \item symmetry factors account for equivalent contractions.
\end{itemize}
Thus the Feynman rules derived earlier from Wick's theorem also emerge directly
from the functional integral.

\section{Connected Generating Functional}

The generating functional \(Z[J]\) includes disconnected diagrams. It is useful
to define the connected generating functional \(W[J]\) by
\[
    Z[J]=e^{iW[J]}.
\]
Equivalently,
\[
    W[J]=-i\ln Z[J].
\]
Functional derivatives of \(W[J]\) generate connected correlation functions.
For example,
\[
    \frac{1}{i}\frac{\delta W[J]}{\delta J(x)}
\]
gives the connected one-point function, and higher derivatives give connected
higher-point functions.

The logarithm removes the disconnected multiplication structure of \(Z[J]\).
This is the functional-integral version of the statement that disconnected
vacuum bubbles exponentiate and cancel in normalized correlators.

\section{The Classical Field}

The source-dependent expectation value of the field is defined by
\[
    \phi_c(x)
    = \frac{\delta W[J]}{\delta J(x)}.
\]
It is called the classical field, although it is really the expectation value of
the quantum field in the presence of a source. When the source is removed,
\(\phi_c\) may vanish or may describe a nontrivial vacuum expectation value.

This quantity is important because it is the natural variable of the effective
action. Instead of describing the theory as a functional of the external source
\(J\), one can describe it as a functional of the expectation value \(\phi_c\).

\section{The Effective Action}

The effective action \(\Gamma[\phi_c]\) is defined by a Legendre transform:
\[
    \Gamma[\phi_c]
    = W[J]-\int\dd^4x\,J(x)\phi_c(x),
\]
where \(J\) is understood as a functional of \(\phi_c\). It satisfies
\[
    \frac{\delta\Gamma[\phi_c]}{\delta\phi_c(x)}=-J(x).
\]
When the source is set to zero, the quantum equation of motion is
\[
    \frac{\delta\Gamma[\phi_c]}{\delta\phi_c(x)}=0.
\]
At tree level, \(\Gamma\) is just the classical action. Loop corrections modify
it. Therefore the effective action is the quantum-corrected action governing
expectation values.

\section{One-Particle-Irreducible Functions}

Functional derivatives of \(\Gamma[\phi_c]\) generate one-particle-irreducible
correlation functions, abbreviated 1PI functions. A diagram is one-particle
irreducible if it cannot be disconnected by cutting a single internal line.

The two-point 1PI function is related to the inverse full propagator. The
four-point 1PI function is related to the quantum-corrected interaction vertex.
This language becomes essential in renormalization, where masses, fields, and
couplings are corrected by loop effects.

The hierarchy is therefore
\[
    Z[J] \longrightarrow \text{all diagrams},
    \qquad
    W[J] \longrightarrow \text{connected diagrams},
    \qquad
    \Gamma[\phi_c] \longrightarrow \text{1PI diagrams}.
\]

\section{Classical Limit and Saddle Points}

The path integral also gives a useful view of the classical limit. When the
action is large compared with \(\hbar\), the integral is dominated by stationary
points of the action:
\[
    \frac{\delta S}{\delta\phi(x)}=0.
\]
This is the classical field equation. Quantum corrections come from fluctuations
around the saddle point.

For a free theory, fluctuations are Gaussian. For an interacting theory,
fluctuations generate vertices and loop diagrams. Thus perturbation theory can
be viewed as a systematic expansion around a saddle point of the action.

\begin{summarybox}
The scalar-field path integral sums over field configurations rather than
particle paths. The source functional \(Z[J]\) generates time-ordered correlation
functions. For a free scalar field, the functional integral is Gaussian and the
inverse quadratic operator is the Feynman propagator. Interactions are generated
by functional derivatives and produce the same diagrammatic expansion as Wick's
theorem. The functionals \(W[J]\) and \(\Gamma[\phi_c]\) organize connected and
one-particle-irreducible diagrams.
\end{summarybox}

\section{Chapter Checklist}

After reading this chapter, one should be able to:
\begin{itemize}
    \item explain the meaning of the functional measure \(\mathcal{D}\phi\);
    \item write the scalar-field generating functional with a source;
    \item derive the free generating functional \(Z_0[J]\);
    \item obtain correlation functions by differentiating with respect to \(J\);
    \item explain why the free path integral is Gaussian;
    \item express interactions as functional derivative operators;
    \item derive the diagrammatic expansion from \(Z[J]\);
    \item define \(W[J]\) and explain connected diagrams;
    \item define the classical field \(\phi_c\);
    \item define the effective action \(\Gamma[\phi_c]\);
    \item explain the meaning of one-particle-irreducible functions.
\end{itemize}
